%============ Chapter1.tex — Introduction ===================
\chapter{Introduction}

\section{Overview}

This report documents the design, implementation, and evaluation of a
research prototype that analyses raw voice recordings for acoustic
patterns associated with Parkinson's Disease (PD) using classical,
explainable Machine Learning (ML). The system covers the complete chain
from an audio file to a cautious screening indication: signal
preprocessing, acoustic feature extraction, model training,
subject-independent validation, external validation on an independent
dataset, and a file-based screening application. This chapter introduces
the medical and engineering motivation behind the project, states the
problem formally, lists the project objectives, defines the scope and
the non-diagnostic framing of the system, and outlines the organization
of the remainder of the report.

\section{Background and Motivation}

Parkinson's Disease is the second most common neurodegenerative
disorder after Alzheimer's disease. It is caused by the progressive loss
of dopamine-producing neurons in regions of the brain that control
movement, and it manifests as tremor, muscular rigidity, slowness of
movement, and reduced movement amplitude. The disease burden is large
and growing: the Global Burden of Disease study estimated that about 6.1
million people were living with PD in 2016, more than double the number
in 1990, a rise driven mainly by population ageing \cite{gbd2018}.

Despite this burden, no simple objective test for PD is available in
routine practice. The diagnosis remains clinical: it is established from
motor signs observed by a neurologist, following criteria such as those
of the Movement Disorder Society (MDS) \cite{postuma2015}. Post-mortem
morphometric evidence indicates that a substantial proportion of the
dopaminergic neurons of the substantia nigra is typically already lost
by the time the characteristic motor signs appear \cite{fearnley1991},
which narrows the window in which future disease-modifying treatments
could be most useful. These facts together motivate the search for
inexpensive, non-invasive measurements that could support earlier
referral of at-risk individuals to clinical evaluation
\cite{postuma2015,fearnley1991}.

The human voice is a promising carrier of such measurements. Speaking is
a motor activity that depends on the coordinated action of the
respiratory muscles, the vocal folds of the larynx, and the
articulators of the vocal tract; the same neuromuscular impairments that
affect limb movement in PD also affect these structures. Speech and
voice impairment --- reduced loudness, monotone pitch, breathiness, and
imprecise articulation, collectively described as hypokinetic
dysarthria --- has been reported in approximately 90\% of a large
clinical sample of PD patients \cite{ho1999}. Importantly for
screening, acoustic abnormalities appear early: quantitative acoustic
analysis detected some form of vocal impairment in about 78\% of
patients with early, untreated PD \cite{rusz2011}, and a longitudinal
case study found measurable reduction of fundamental-frequency
variability in recordings made about five years before the clinical
diagnosis \cite{harel2004}.

A voice recording can be captured in seconds with a commodity
microphone or smartphone, at negligible cost, with no discomfort, and
repeatedly. If acoustic measurements of such recordings can be related
to PD status with useful reliability, voice analysis could serve as a
low-cost screening-support signal, flagging individuals for whom a
clinical examination is advisable. This prospect has produced a large
research literature and several public datasets, including the Mobile
Device Voice Recordings at King's College London (MDVR-KCL) corpus used
in this project \cite{jaeger2019}.

Two considerations shaped the engineering approach taken here. First,
in a biomedical setting, the ability to explain \emph{why} a model
produced its output is valued alongside raw performance; classical ML
models operating on physiologically interpretable acoustic features
were therefore preferred over deep learning, whose learned
representations are difficult to interpret and whose data requirements
far exceed the 37 subjects available. Second, published results in this
field frequently report accuracies above 95\%, which, as examined in
Chapters 2 and 4, often stems from evaluation protocols that let the
same speaker appear in both training and test data. Producing an
honest, subject-independent, and externally tested estimate of
performance was therefore treated not as an obstacle but as a central
contribution of the project.

\section{Problem Statement}

The problem addressed by this project is stated as follows: given a raw
voice recording of continuous speech, determine --- by means of
reproducible signal processing and classical machine learning ---
whether the acoustic pattern of the recording is closer to that of a
group of speakers with PD or to that of a group of healthy control
speakers, and report the result with an honest, subject-independent
estimate of reliability and with wording appropriate to a
non-diagnostic research tool.

Solving this problem requires answering several subsidiary questions:
which acoustic measurements meaningfully capture Parkinson's-related
voice changes in continuous speech; how such measurements must be
extracted so that training, evaluation, and prediction remain strictly
consistent; how a model must be validated when several recordings
originate from the same person; how much of the measured performance
survives transfer to a different dataset, language, and recording
channel; and how the output of such a system must be presented to avoid
overstating what it knows.

\section{Project Objectives}

The objectives were divided into a primary set, which defines the
minimum complete system, and an extended set, which was undertaken after
the primary objectives were met.

\subsection{Primary Objectives}

\begin{enumerate}
  \item Inspect and verify the MDVR-KCL dataset before any modelling:
        file integrity, class labels, subject identities, and the
        feasibility of subject-independent validation.
  \item Implement a consistent preprocessing pipeline for raw audio
        (mono conversion, fixed sample rate, offset removal, silence
        trimming, and amplitude normalization) that deliberately avoids
        operations that would distort clinically meaningful signal
        irregularities.
  \item Implement, as a core contribution, the extraction of
        physiologically motivated acoustic features: fundamental
        frequency statistics, jitter, shimmer, harmonics-to-noise
        ratio, and mel-frequency cepstral coefficient statistics.
  \item Train and compare classical ML classifiers (logistic
        regression, support vector machine, random forest, and
        multilayer perceptron) against a majority-class baseline.
  \item Evaluate all models with strictly subject-independent
        cross-validation, reporting balanced accuracy, sensitivity,
        specificity, F1 score, confusion matrices, and the area under
        the receiver operating characteristic curve, together with
        their variability across folds and random seeds.
  \item Build a file-based screening prototype (browser application and
        command-line interface) that reuses the identical pipeline and
        presents results with mandatory non-diagnostic wording.
\end{enumerate}

\subsection{Extended Objectives}

\begin{enumerate}
  \setcounter{enumi}{6}
  \item Improve classification performance through recording
        segmentation and additional continuous-speech features
        (smoothed cepstral peak prominence, pause statistics, and
        delta mel-frequency cepstral coefficients) under
        model-selection rules declared before experimentation, so that
        the improvement process itself cannot overfit the validation
        data.
  \item Demonstrate quantitatively, on the project's own data, how an
        incorrect validation design inflates reported performance.
  \item Validate the final, frozen model externally on an independent
        dataset recorded in a different language and under different
        conditions, and report the resulting generalization gap.
  \item Extend the prototype with uncertainty-aware reporting: an
        inconclusive score band, recording-condition warnings, and a
        microphone demonstration mode explicitly framed as
        out-of-domain.
\end{enumerate}

All ten objectives were achieved; the evidence is presented in
Chapter~4.

\section{Scope and Non-Diagnostic Framing}

The system developed in this project is a research screening-support
prototype. Throughout this report, and inside the software itself, its
output is described as a non-diagnostic indication of similarity
between the acoustic pattern of a recording and the patterns of the
dataset groups. The system does not diagnose Parkinson's Disease, does
not measure any person's medical risk, and must not replace evaluation
by a qualified healthcare professional. Many conditions other than PD
--- among them common respiratory infections, ageing, smoking, and
fatigue --- also alter voice quality, and the datasets used are far too
small and too homogeneous to support any clinical claim. The
quantitative results presented in this report characterize performance
over the datasets studied, not the accuracy of any statement about an
individual. A dedicated discussion of limitations and the ethical
constraints applied to the system's wording is given in
Section~4.12.

Within these boundaries, the scope of the work comprises: the MDVR-KCL
dataset as the training corpus; raw audio as the only input
representation, with all features extracted by the project's own code;
classical machine learning as the modelling approach; an independent
Italian corpus as an external test set; and a file-based prototype with
an additional microphone demonstration mode. Live clinical deployment,
longitudinal patient tracking, and deep learning approaches are outside
the scope of this project.

\section{Report Organization}

The remainder of this report is organized as follows.

\textbf{Chapter 2 --- Background and Literature Review} presents the
clinical background of PD and its effect on speech production, defines
the acoustic measurements used in voice analysis, introduces the
classical ML models employed, examines the data-leakage problem that
affects much of the published literature, and reviews related work on
the datasets used here.

\textbf{Chapter 3 --- Methodology} specifies the complete system: the
dataset inspection protocol, the preprocessing chain, recording
segmentation, the 74-feature extraction stage, the classification
pipeline, the subject-independent validation protocol, the pre-declared
model-selection rules, the external validation design, and the
prototype.

\textbf{Chapter 4 --- Results and Discussion} reports the outcomes of
every stage: dataset inspection, feature validation, baseline
modelling, the model-improvement experiments, the final model, feature
importance, the validation-design demonstration, external validation,
and the prototype, followed by a general discussion and a dedicated
section on limitations and ethical considerations.

\textbf{Chapter 5 --- Conclusion and Future Work} summarizes the
findings in points and lists recommendations for future development.

Four appendices provide the complete feature inventory, the full
experimental results, a reproduction guide for the software repository,
and key code excerpts.
